%! AP Statistics — Unit 1, Lesson 1.5 — Slides (Beamer)
\documentclass[aspectratio=169, 11pt]{beamer}
\usepackage{apstats-beamer}

%=============================================================================
\begin{document}

% ─────────────────────────────────────────────────────────────────────────────
% SLIDE 1 — LESSON TITLE
% ─────────────────────────────────────────────────────────────────────────────
{
\setbeamercolor{background canvas}{bg=navy}
\begin{frame}[plain]
  \begin{tikzpicture}[remember picture, overlay]
    \fill[navylight] (current page.north west)
      rectangle ([yshift=-1.35cm]current page.north east);
  \end{tikzpicture}
  \vspace{2.8cm}
  \hspace{0.55cm}%
  \begin{minipage}{0.9\textwidth}
    {\color{goldacc}\bfseries\small LESSON 1.5}\\[0.5cm]
    {\color{white}\bfseries\LARGE Representing a Quantitative Variable\\[0.25cm]
      with Graphs}\\[1.4cm]
    {\color{skymid}\small AP Statistics~$\cdot$~Unit 1: Exploring One-Variable Data}
  \end{minipage}
\end{frame}
}

% ─────────────────────────────────────────────────────────────────────────────
% SLIDE 2 — WARM-UP
% ─────────────────────────────────────────────────────────────────────────────
{
\setbeamercolor{background canvas}{bg=sky}
\begin{frame}[t, plain]
  \navyheader{Warm-Up}
  \hspace{0.55cm}%
  \begin{minipage}{0.9\textwidth}
    \vspace{0.3em}
    \small
    \textit{Look at the dotplot on your warm-up sheet showing books read this summer.}
    \begin{enumerate}
      \setlength{\itemsep}{0.6em}\setlength{\topsep}{0em}
      \item What is the variable? Is it quantitative or categorical?
      \item What value appears most often? How many students chose it?
      \item What is the range of the distribution?
      \item How is a dotplot different from a bar chart?
    \end{enumerate}
    \vspace{0.4cm}
    {\color{navy}\itshape\small
      Today we learn three displays for quantitative data and practice
      reading each one.}
  \end{minipage}
\end{frame}
}

% ─────────────────────────────────────────────────────────────────────────────
% SLIDE 3 — PRIMARY OBJECTIVE
% ─────────────────────────────────────────────────────────────────────────────
{
\setbeamercolor{background canvas}{bg=sky}
\begin{frame}[t, plain]
  \begin{tikzpicture}[remember picture, overlay]
    \fill[navy] (current page.north west)
      rectangle ([yshift=-0.25cm]current page.north east);
  \end{tikzpicture}
  \vspace{0.45cm}\par
  \hspace{0.55cm}%
  \begin{minipage}{0.9\textwidth}

    \sectionlabel{Primary Objective}

    \normalsize
    Students will \textbf{read and interpret} dotplots, stemplots, and histograms
    for quantitative data, and will identify which display is most appropriate for
    a given dataset.

    \vspace{0.5cm}
    \textbf{Learning Targets:}
    \begin{enumerate}\small
      \setlength{\itemsep}{0.4em}
      \item Read a dotplot: identify center, spread, shape, gaps, clusters.
      \item Complete and interpret a stemplot; read a back-to-back stemplot.
      \item Interpret a histogram from a frequency table; explain why bars touch.
      \item Choose the right display based on dataset size and purpose.
    \end{enumerate}
  \end{minipage}
\end{frame}
}

% ─────────────────────────────────────────────────────────────────────────────
% SLIDE 4 — HOOK: THREE DISPLAYS, ONE DATASET
% ─────────────────────────────────────────────────────────────────────────────
{
\setbeamercolor{background canvas}{bg=goldbg}
\begin{frame}[t, plain]
  \navyheader{Hook: Three Displays, One Dataset}
  \hspace{0.55cm}%
  \begin{minipage}{0.9\textwidth}
    \vspace{0.2em}\small
    Ages of 15 summer camp participants:
    \[10,10,10,10,10,11,11,11,11,12,12,12,13,13,14\]

    \vspace{0.3cm}
    \begin{tabular}{p{0.28\textwidth} p{0.28\textwidth} p{0.28\textwidth}}
    \textbf{Dotplot} & \textbf{Stemplot} & \textbf{Histogram} \\
    \hline
    Every value visible & Sorted; values preserved & Shape visible; values grouped \\
    Best for small $n$ & Best for moderate $n$ & Best for large $n$ \\
    \end{tabular}

    \vspace{0.5cm}
    {\color{navy}\bfseries\small Discussion:} What does each display make easy to see?
    What does each one hide?
  \end{minipage}
\end{frame}
}

% ─────────────────────────────────────────────────────────────────────────────
% SLIDE 5 — DOTPLOTS
% ─────────────────────────────────────────────────────────────────────────────
{
\setbeamercolor{background canvas}{bg=white}
\begin{frame}[t, plain]
  \navyheader{Display 1: Dotplots}
  \hspace{0.55cm}%
  \begin{minipage}{0.9\textwidth}
    \vspace{0.2em}
    \begin{columns}[t]
      \column{0.45\textwidth}
        \small
        \textbf{\color{navy}How to read:}
        \begin{itemize}\setlength{\itemsep}{0.4em}
          \item Horizontal axis = values of the variable
          \item Each dot = one individual
          \item Stack height = frequency
          \item Look for: center, spread, shape, gaps, clusters
        \end{itemize}

      \column{0.45\textwidth}
        \small
        \textbf{\color{navy}Best when:}
        \begin{itemize}\setlength{\itemsep}{0.4em}
          \item $n$ is small (roughly $\le 30$--50)
          \item You want to see every individual value
        \end{itemize}

        \vspace{0.3cm}
        \textbf{\color{navy}Limitation:}\\
        Impractical for large $n$ --- dots pile up
    \end{columns}

    \vspace{0.5cm}
    {\color{navy}\bfseries\small Key point:}
    Dotplot dots \textit{do} stack touching because the number line is continuous
    --- just like histogram bars.
  \end{minipage}
\end{frame}
}

% ─────────────────────────────────────────────────────────────────────────────
% SLIDE 6 — STEMPLOTS
% ─────────────────────────────────────────────────────────────────────────────
{
\setbeamercolor{background canvas}{bg=white}
\begin{frame}[t, plain]
  \navyheader{Display 2: Stemplots (Stem-and-Leaf)}
  \hspace{0.55cm}%
  \begin{minipage}{0.9\textwidth}
    \vspace{0.2em}
    \begin{columns}[t]
      \column{0.45\textwidth}
        \small
        \textbf{\color{navy}Structure:}
        \begin{itemize}\setlength{\itemsep}{0.4em}
          \item Stem = leading digit(s)
          \item Leaf = final digit
          \item Leaves in ascending order
          \item Always write a key: \textbf{7 $|$ 4 = 74}
        \end{itemize}

        \vspace{0.3cm}
        \textbf{\color{navy}Advantage:}\\
        Data is automatically sorted $\Rightarrow$ median is easy to find

      \column{0.45\textwidth}
        \small
        \textbf{\color{navy}Example:} Quiz scores\\
        54, 58, 61, 65, 67, 72, 78, 79, 85

        \vspace{0.2cm}
        \begin{tabular}{r|l}
        \textbf{Stem} & \textbf{Leaf} \\
        \hline
        5 & 4\; 8 \\
        6 & 1\; 5\; 7 \\
        7 & 2\; 8\; 9 \\
        8 & 5 \\
        \end{tabular}

        \vspace{0.2cm}
        \small Median? \hspace{0.5cm} $n = 9$ $\Rightarrow$ 5th value $= 67$
    \end{columns}
  \end{minipage}
\end{frame}
}

% ─────────────────────────────────────────────────────────────────────────────
% SLIDE 7 — HISTOGRAMS
% ─────────────────────────────────────────────────────────────────────────────
{
\setbeamercolor{background canvas}{bg=white}
\begin{frame}[t, plain]
  \navyheader{Display 3: Histograms}
  \hspace{0.55cm}%
  \begin{minipage}{0.9\textwidth}
    \vspace{0.2em}
    \begin{columns}[t]
      \column{0.48\textwidth}
        \small
        \textbf{\color{navy}Key features:}
        \begin{itemize}\setlength{\itemsep}{0.35em}
          \item Equal-width bins (class intervals)
          \item Bar height = frequency (or relative frequency)
          \item \textbf{Bars touch} --- adjacent bins share a boundary
          \item Vertical axis starts at 0
        \end{itemize}

        \vspace{0.3cm}
        \textbf{\color{navy}Best when:}\\
        $n$ is large ($> 50$); shape is more important than individual values

      \column{0.45\textwidth}
        \small
        \textbf{\color{navy}Histogram vs.\ bar chart:}

        \vspace{0.2cm}
        \begin{tabular}{ll}
        Bars \textbf{touch} & $\Rightarrow$ quantitative data \\
        Bars \textbf{separate} & $\Rightarrow$ categorical data \\
        Order matters & $\Rightarrow$ histogram \\
        Order arbitrary & $\Rightarrow$ bar chart \\
        \end{tabular}

        \vspace{0.3cm}
        \textbf{\color{redacc}Common error:}\\
        Drawing gaps between histogram bars (it's not a bar chart!)
    \end{columns}
  \end{minipage}
\end{frame}
}

% ─────────────────────────────────────────────────────────────────────────────
% SLIDE 8 — CHOOSING A DISPLAY
% ─────────────────────────────────────────────────────────────────────────────
{
\setbeamercolor{background canvas}{bg=greenbg}
\begin{frame}[t, plain]
  \navyheader{Choosing the Right Display}
  \hspace{0.55cm}%
  \begin{minipage}{0.9\textwidth}
    \vspace{0.3em}\small

    \begin{tabular}{p{0.28\textwidth} p{0.25\textwidth} p{0.3\textwidth}}
    \rowcolor{sky}
    \textbf{Display} & \textbf{Best for} & \textbf{Trade-off} \\
    \hline
    Dotplot & Small $n$ ($\le 30$--50) & Impractical at large $n$ \\[4pt]
    Stemplot & Moderate $n$ (15--50) & Awkward with large values or non-integers \\[4pt]
    Histogram & Large $n$ ($> 50$) & Loses individual values; bin width matters \\
    \end{tabular}

    \vspace{0.5cm}
    \textbf{\color{navy}Quick rules of thumb:}
    \begin{itemize}\setlength{\itemsep}{0.4em}
      \item Comparing two groups of 10--30? $\Rightarrow$ Back-to-back stemplot
      \item Dataset with 500 observations? $\Rightarrow$ Histogram
      \item 8 runners' finishing times? $\Rightarrow$ Dotplot
    \end{itemize}
  \end{minipage}
\end{frame}
}

% ─────────────────────────────────────────────────────────────────────────────
% SLIDE 9 — GROUP ACTIVITY LAUNCH
% ─────────────────────────────────────────────────────────────────────────────
{
\setbeamercolor{background canvas}{bg=sky}
\begin{frame}[t, plain]
  \navyheader{Group Activity}
  \hspace{0.55cm}%
  \begin{minipage}{0.9\textwidth}
    \vspace{0.3em}\small
    \textit{Work through your activity sheet in your group.}

    \vspace{0.4cm}
    \begin{enumerate}
      \setlength{\itemsep}{0.5em}
      \item \textbf{Part 1 (8 min):} Read and analyze the provided dotplot of
            seedling heights.
      \item \textbf{Part 2 (8 min):} Complete the back-to-back stemplot for two
            AP Statistics classes and compare their distributions.
      \item \textbf{Part 3 (5 min):} Complete a frequency table and answer
            interpretation questions about a histogram.
      \item \textbf{Part 4 (3 min):} Individual reflection --- the most important
            trade-off when choosing a display.
    \end{enumerate}

    \vspace{0.4cm}
    {\color{navy}\itshape\small
      Be prepared to share your Part 2 stemplot and one observation from Part 3.}
  \end{minipage}
\end{frame}
}

% ─────────────────────────────────────────────────────────────────────────────
% SLIDE 10 — EXIT TICKET
% ─────────────────────────────────────────────────────────────────────────────
{
\setbeamercolor{background canvas}{bg=redbg}
\begin{frame}[t, plain]
  \navyheader{Exit Ticket}
  \hspace{0.55cm}%
  \begin{minipage}{0.9\textwidth}
    \vspace{0.3em}\small
    \textit{Complete the exit ticket independently (5 minutes).}

    \vspace{0.3cm}
    A stem-and-leaf plot shows resting pulse rates for 12 patients:

    \vspace{0.2cm}
    \begin{tabular}{r|l}
    \textbf{Stem} & \textbf{Leaf} \\
    \hline
    6 & 2\; 4\; 5\; 8 \\
    7 & 0\; 2\; 2\; 4\; 5\; 6\; 8 \\
    8 & 0 \\
    \end{tabular}
    \quad Key: \textbf{7 $|$ 2} = 72 bpm

    \vspace{0.3cm}
    \begin{enumerate}\setlength{\itemsep}{0.35em}
      \item How many patients have a pulse rate in the 70s?
      \item What is the range?
      \item What is the median?
      \item Describe the shape.
      \item Write one AP-style sentence describing the distribution.
    \end{enumerate}
  \end{minipage}
\end{frame}
}

% ─────────────────────────────────────────────────────────────────────────────
% SLIDE 11 — LESSON WRAP-UP AND PREVIEW
% ─────────────────────────────────────────────────────────────────────────────
{
\setbeamercolor{background canvas}{bg=navy}
\begin{frame}[plain]
  \vspace{1.8cm}
  \hspace{0.55cm}%
  \begin{minipage}{0.9\textwidth}
    {\color{goldacc}\bfseries\large Wrap-Up}\\[0.5cm]
    {\color{white}\normalsize
    Today you learned to read three displays for quantitative data:
    dotplots, stemplots, and histograms.}\\[0.4cm]
    {\color{skymid}\normalsize
    \textbf{Coming next --- Lesson 1.6:}\\
    Use these graphs to describe a distribution with the
    \textbf{SOCS} framework: \textbf{S}hape, \textbf{O}utliers,
    \textbf{C}enter, \textbf{S}pread.}\\[0.4cm]
    {\color{white!80!black}\small
    Before next class: complete the homework and review today's graphs.
    Practice describing what you see in 1--2 complete sentences.}
  \end{minipage}
\end{frame}
}

\end{document}
